% !TeX root = ../papers/nc-p3.tex
\section{二平面分解の正則特殊化と追加単純成分の全点}
\label{sec:split-regular-limits}
\begin{theorem}\label{thm:split-regular-limits}
分数体上の窓が可逆二平面分解を持つか、または可逆二平面分解のある閉軌道へ流入するとする。
任意の基底変更後に標準正則な特殊点を持てば、その特殊点は半安定で、
商で対応する唯一の閉軌道も可逆二平面分解を持ち標準正則である。
従ってこのような閉軌道への全流入の合併は正則軌跡で特殊化に閉じる。
\end{theorem}
\begin{proof}[整模型の体積と四つの純成分]
全体テンソル $w$ の $Y_iY_{i+1}X_{i+2}X_{i+3}$ 成分を
$\kappa_i s_i\otimes r_{i+2}$ と書く。これを色が混合した二組の間で切った四次正方行列を $C_i$ とする。
その行列式は符号を除き $\kappa_i^4(\det s_i\det r_{i+2})^2$ である。
有限拡大後、$X_i$ の各位置の高位文字を選ぶ際、四つの $r_i$ で
高位・高位の係数を同時に零にできる。実際、双対射影直線上の高位係数を読む点を、
各非退化形式の直交で次の点へ送り、四つの射影同型の合成の不動点から始めればよい。
$Y_i$ でも同様に選ぶ。二色を $X<Y$ とし、$r_i,s_i$ の高位・低位の係数を $r_i',s_i'$ と書く。
どちらも非零である。全色順の降順対が先頭語のPBW表示となる。まず色だけを並べ替え、
次に純二文字の降順対を昇順にする。前者は $O(n^2)$ 回であり、後者で $10$ を $01$ に
するたび逆転数が減り、$00$ にするたび高位文字数が減って逆転数の増加は高々 $n$ である。
よって各項の書換え経路が $O(n^2)$ で、各係数の付値に定理\ref{thm:pbw-valuative-boundaries}の二次下界がある。

四つの巡回における最上降順語の係数は $p_i=\kappa_i s_i' r_{i+2}'$。
モニックな最上元での $1+3$ 行列の行列式は符号を除き
\[
 \frac{\kappa_i^2\kappa_{i+1}^2\det(s_i)\det(r_{i+3})(r_{i+2}')^2(s_{i+1}')^2}{p_i^4}.
\]
実際、先頭が $Y$ の二行では残りが $YXX$、先頭が $X$ の二行では $YYX$ であり、
各 $2\times2$ ブロックを上記純成分の降順対で縮約した式である。
四行列式の付値和 $\delta$ は $\prod_i\det(s_i)\det(r_i)/(s_i'r_i')^2$ の付値となる。

一般の正則格子の座標への基底変更 $g$ を有限拡大と四位置のスカラー倍により
各位置特殊線形に取る。$gw$ の最小係数付値を $\ell$ とする。全巡回で同じであり、
正則側の最上格子は原始的なので、参照モニック基底との比較付値和は
$M=4\ell-\sum_i\operatorname{ord}p_i$ である。
定理\ref{thm:pbw-valuative-boundaries}の証明中の $V=M-\delta/2$ に $V=0$ を入れれば
\[
              16\ell=\sum_{i=0}^3\operatorname{ord}\det C_i .
\]
特にこれは純関係代表の倍率や高位文字の選び方に依らない。
\end{proof}
\begin{proof}[閉軌道への流入を一般ファイバーに許す場合]
一般テンソルを $z$ とし、一般体の有限拡大と位置別の基底変更後に、
$\lambda(u)z\to w$ ($u\to0$) とする。ここで $u$ は元の離散付値の一様化元とは独立であり、
$\lambda$ は特殊線形群の余指標である。半安定商の閉軌道への一径数判定によりこの形を使える。
$w$ を保つ $\lambda$ と可逆分解の余指標は、後者の印を $w$ の安定群で共役すれば
同じ極大トーラスに入る。このため $X,Y$ を $\lambda$ 不変に取れる。

二色各々について、純二文字のモニック書換えに負付値係数を生じないよう、
さらに $\lambda$ 固有基底を選べる。実際、非退化純形式の直交が定める射影直線同型に沿って
基点の基底を順に運べば、前三辺の形式はスカラー倍を除き交代形式になり、最後の一辺だけに
四同型の合成 $B\in\mathrm{PGL}_2$ が残る。この輸送は $\lambda$ の射影作用と可換である。
基点で作用が非自明なら $B$ はそのトーラスと可換なので同じ二固有線を持つ。
作用がスカラーなら自由にJordan基底を取れる。半単純な場合、高位と低位を交換して
最後のモニック規則 $10\to c\,01$ で $\operatorname{ord} c\ge0$ に取る。
非半単純なら作用はスカラーで、通常のJordan行列の代表を選んで
$10\to c\,01+d\,00$ の両係数を整にできる。スカラー $B$ も含む。
こうして純書換えの回数は付値下界に影響しない。

$z$ の関係を降順対の六列で行消去する。その小行列式の非消滅は $w$ での
非消滅と $\lambda$ 流入から従う。Grassmannチャートのグラフ係数は固有座標だから、
各降順対の右辺で同じ余重みの部分は $w$ の規則、残りは全て高い余重みを持つ
（行が張る関係部分空間への作用で右辺と左辺の文字重み差を取る）。
低い余重みを優先して同重みでは元の左辞書順にすると全規則は降下する。
$z$ は正則背景への良い還元で標準正則なので次元比較から昇順語は全次数で基底である。
最上元の各巡回の降順語は $w$ に非零係数を持つので、$z$ でも同じ $p_i$ である。
また $1+3$ 行列式では、背景の非零な置換項があるから、各置換項に現れる
四語の余重みの総和は背景重みの四倍である（行文字と列三文字の多重集合が一定）。
背景より高い項を一つでも含めば背景より低い項も要り、後者は存在しない。
従って四つのモニック行列式も全て変わらない。

長さ $n$ の各書換え経路上、余重みが上がる回数は、全語の余重みの幅から $O(n)$ である。
一回の増加で二色の逆転数は高々 $2n$ 増える。同重みで異色を動かせば逆転数が1減り、
同色の背景書換えでは不変である。従って同重みの異色交換も計 $O(n^2)$ 回に限られる。
背景の同色係数は全て整で、それ以外の各一歩の係数付値は一定値以上だから、
全語の正規形係数に $-Cn^2$ の下界を得る。定理\ref{thm:pbw-valuative-boundaries}の
トーションの議論はこの余重み優先順でも成立する。中央の微分の初期形はやはり
降順頭部と昇順尾部の単位写像であり、固定長における順序は有限で連接と両立するからである。
ゆえに $z$ の正則特殊模型への特殊線形比較での最小付値 $\ell$ は
先の体積公式の右辺（$w$ で計算）で決まる。
$z$ と $w$ の斉次特殊線形不変式は最初の余指標極限で等しいため、
次の整従属性の議論はこの同じ $\ell$ で適用できる。
\end{proof}
\begin{proof}[分解を保つ群の商への延長]
$G=\prod_i\mathrm{SL}_4$、各位置の $X,Y$ に重み $-1,1$ を与えるトーラスを $T$、
その中心化群を $C$ とする。周囲の四テンソル表現を $V$ とし、$V^T$ 上の
$C$ 不変式環は $k[V]^G$ の制限上整であることを確認する。
零錐を比較すれば十分である。実際、$C$ 軌道閉包の閉軌道において、複素数上で
最大コンパクト群に関する極小ベクトルを取ると、各位置の一文字Gramは $T$ 不変性で
二色間の非対角成分が零である。中心化群のモーメントが零なので各位置のGram全体がスカラーとなり、
全群でも極小ベクトルで閉軌道である。従って $G$ 零錐の制限は $C$ 零錐に等しい。
この代数的な結論は標数零代数閉体へ基底拡大で移る。
有限生成された二つの次数付き不変式環で、正次数制限のイデアルの根基は
$C$ 商の頂点イデアルだから、その剰余は有限次元である。
斉次な剰余基底を持ち上げ次数帰納すれば、制限環上の加群有限性を得る。

$\pi^{-\ell}w$ の $G$ 不変式は、正則特殊模型における原始 $z$（分解を初めから持つ場合は原始 $w$）の係数から整である。整従属性により全 $C$ 不変式も整であり、
その剰余が $V^T/\!/C$ の一点を定める。そのファイバーの閉軌道代表を $v$ とする。
$\det C_0\det C_2,\det C_1\det C_3$ はいずれも次数8の $C$ 不変式である。
各位置で一方の因子に $X$、他方に $Y$ の行列式が同じ指数2で掛かるからである。
両不変式は整で、その積は体積公式から単元だから $v$ でも四つの $C_i$ は全て可逆である。
全二切断が階数6以下、四純成分が純対どうしの切断で階数1以下という閉条件も $v$ に残る。
実際、これらは $C$ 不変な閉集合でその良い商への像は閉で、像上の閉軌道はこの集合にある。
従って四純成分は全て非退化二次行列の積であり、両隣接切断は色別に少なくとも $(1,4,1)$、
一文字も全階数である。定理\ref{thm:split-inverse-regularity}で $v$ は可逆分解の正則窓となる。
前段の極小ベクトル論法で $v$ は $G$ についても閉軌道である。
正則模型の原始特殊テンソルは $v$ と全 $G$ 不変式が等しく、
$v\ne0$ により半安定で、同じ閉軌道に対応する。
\end{proof}

\begin{theorem}\label{thm:extra-simple-all-regular}
定理\ref{thm:split-simple-surface}の射影60次元の追加単純成分の
標準正則点は、$b\in k\setminus\{0,-1\}$ の表示済みの各一類と、
$b=-1$ における片側反転の閉軌道一類・その例外二鎖Jordan一類で尽きる。
$b=0,\infty$ に正則な特殊化はない。
\end{theorem}
\begin{proof}
追加表示 $(1,-\iota t,\iota,t)$, $L=ED_0$ を用いる。
$\infty$ に向かう $t$ は逆数で $0$ にでき、$t\to\pm1$ も符号と
一添字移動の $t\mapsto(t-1)/(t+1)$ および逆数で $0$ に移る。
従って $t$ の付値 $e>0$ を扱えばよい。四位置の文字の標識を順に
$(0,1,2,3),(1,0,2,3),(0,1,3,2),(1,0,3,2)$ とする。
モニック $1+3$ 行列式の積は1で、四つの降順語係数は $((1-t^2)/(2t),1,1,1)$ となる。
従って正則な特殊点があれば上の特殊線形規格化で $\ell=-e/4$ である。
ところが反対切断 $02|13$ の行列式は次数16の全群不変式で
\[
 \det w_{02|13}=
 \frac{(t^2+1)^4(t^2-2t-1)^2(t^2+2t-1)^2}{256t^8}.
\]
\path{computations/verify_split_extra_boundary.py} は四位置の関係の先頭語、体積と
この恒等式を有理関数体で検算する。原始模型に直した不変式の付値は $-8e-16\ell=-4e<0$ で矛盾する。

既約成分の任意の点を稠密な表示からの曲線の特殊点とし、必要なら分数体を有限拡大する。
残る $t$ は単元で剰余は $\pm1$ でなく、$b=f(t)^2$ が $-1$ なら
定理\ref{thm:mixed-flip-slice}の平面で正則な閉模型に延びる。他は元の単純表示で正則かつ閉軌道に延びる。
二正則格子の比較によりその商ファイバーに限られる。
$b=-1$ の二類は同平面で実現される。他の閉軌道の射影次元は59で、
それを境界に持つ別軌道が成分にあれば次元60で稠密となり、
成分上の非定数不変有理関数 $b$ と矛盾する。
\end{proof}
